\subsection{ARCH-136-DEGREE-TWO-RESIDUE: the exact coefficient obstruction}
\label{sec:136degree2residue}

\paragraph{Sheet and scope.}
This gives the precise remaining relative obstruction in the construction
ARCH-136-DEGREE-TWO-FREE. The admissibility bounds are those of cote136,
physical pp.~117 and119, footers116 and118, in
\texttt{ARCHIVE\_136\_scan117.tex} and
\texttt{ARCHIVE\_136\_scan119.tex}; the common-source question is on
physical p.~121, footer120, \texttt{ARCHIVE\_136\_scan121.tex}.
All algebras and maps remain strict differential divided-power
$S=\Gamma_{\mathbf Z}(T)$-algebras. The proof is elementary, with tag
none beyond the already specified admissibility and Dold--Kan inputs.

\paragraph{Statement.}
Let $C_1,\ldots,C_t$ be admissible theories with $(C_j)_0=S$,
and let $f_j:A\to C_j$ be strict maps with $A_0=S$.
There is a single weight bound $W$, independent of $n$ and of the
positive integer $N$, for the following assertions.
For $w\geq W$, fix
\[
 z\in A_n^{w,2},\qquad b\in A_n^{w,1},\qquad
                    dz=0,\quad db=Nz.                 \tag{136.DR1}
\]
Write $V_n$ for the vertices of $\Delta[n]$ and put
\[
 Q_{n,w,N}=(S^w)^{n+1}/
              \bigl(\operatorname{diag}S^w+N(S^w)^{n+1}\bigr).
                                                               \tag{136.DR2}
\]
In each target choose $c_j,t_j$ with
\[
 dc_j=f_j(z),\qquad dt_j=Nc_j-f_j(b),
              \qquad |c_j|=(w,1),\quad |t_j|=(w,0).
                                                               \tag{136.DR3}
\]
Such choices exist, and
\[
                 \kappa_j(z,b,N)=[-t_j|_{V_n}]
                           \in Q_{n,w,N}                \tag{136.DR4}
\]
is independent of both choices. A free simplicial Tate-pair
attachment under $A$, adjoining $c,t$ with
$dc=z$, $dt=Nc-b$, can extend every $f_j$ and preserve $A_0=S$
by prescribing common coefficient values for $t$ on vertices
if and only if all the classes (136.DR4) are equal.
The attachment then remains rationally a free unit-disk extension
in each simplex. This criterion is for that prescribed attachment
and those prescribed maps; it is not a criterion for connectedness
of the original category.

Suppose additionally that, in the weight under consideration,
$H^0(A_n)=S^w$ and $H^1(A_n)=0$ for every $n$, and that the
form-zero row of $A$ is acyclic. Then for every pair of targets
there is a canonical homomorphism
\[
 \delta_{jk,n,w,N}:H^2(A_n^{w,*})[N]\longrightarrow Q_{n,w,N}
                                                               \tag{136.DR5}
\]
with
\[
 \delta_{jk,n,w,N}([z])=
                  \kappa_j(z,b,N)-\kappa_k(z,b,N).
                                                               \tag{136.DR6}
\]
In particular the difference is independent of the chosen relation
$db=Nz$. These homomorphisms are the exact compatibility obstruction
for the fixed-degree source of ARCH-136-DEGREE-TWO-FREE.

\begin{proof}
Choose $W\geq1$ so that all targets have coefficient $H^0$ and
zero $H^1,H^2$ in every simplex at weights at least $W$, and so
their form-zero rows are acyclic there. The finite maximum of
these bounds has no dependence on $n$ or $N$. Acyclicity of a row
implies that a simplicial module map prescribed on the vertices
of $\Delta[n]$ extends to the whole simplex: the relative
normalized free chains are extended one basis element at a time,
using acyclicity to fill its already assigned boundary. Thus
\[
                 C_{j,n}^{w,0}\longrightarrow(S^w)^{n+1}
                                                               \tag{136.DR7}
\]
is surjective. This argument uses no chosen linear contraction.

The two vertical vanishings give (136.DR3). If $c'_j,t'_j$ are
another choice, then $c'_j-c_j=du$ for some zero-form $u$, and
$t'_j-t_j-Nu$ is closed in degree zero. It is therefore one
coefficient in $S^w$, constant on all vertices. The changes of
vertex values are exactly a multiple of $N$ plus a diagonal
coefficient. This proves independence in (136.DR4).

Necessity of equality follows because the image of a single
source $t$ has the same vertex coefficients in every target.
Conversely, let $v_j=t_j|_{V_n}$ and choose $v=v_1$.
Equality of the residue classes says that for each $j$ there
are a vector $u_j\in(S^w)^{n+1}$ and a coefficient $a_j\in S^w$
such that
\[
                         v_j+Nu_j+\operatorname{diag}(a_j)=v.
\]
Lift $u_j$ to a zero-form $\widetilde u_j$ by (136.DR7), and
replace the chosen pair by
\[
 c_j+d\widetilde u_j,\qquad
                  t_j+N\widetilde u_j+a_j.
\]
The differential identities are unchanged, and all vertex
values of the new $t_j$ are $v$. The $c_j$ have zero vertex
values because the targets have no positive forms in simplex zero.

Adjoin generators parametrized by $\mathbf Z[\Delta[n]]$,
imposing on the submodule $\mathbf Z[V_n]$ the relations
$c=0$ and $t=v$. Give them the differentials in the statement.
The old $z,b$ vanish on vertices, so every imposed relation
is stable under the differential. The PD ideal of the relations
is stable by the special derivation identity. This constructs
an actual simplicial strict PD extension. The adjusted target
pairs define its maps by the universal property. In every
simplex the inclusion of the vertex module has a free complement,
so the underlying extension is free on the remaining generators.
There are no new generators in simplex zero. Over $\mathbf Q$
the change $c'=c-b/N$ gives $dc'=0$, $dt=Nc'$, an even unit
disk, in each such free complement. Its integral divided-power
Poincare contraction after scaling by $N$ proves rational
acyclicity of the extension. This proves the criterion.

Here is also a cohomological description which proves the last
assertions and makes the obstruction canonical. For any one of
these algebras $D$, set
\[
 K_D=\ker\bigl(D_n^{w,*}\longrightarrow(S^w)^{n+1}[0]\bigr).
\]
The displayed restriction is degreewise surjective by (136.DR7),
or by the same argument for $A$. The exact sequence of complexes
and the assumed degree-zero and degree-one cohomology give
\[
 H^0(K_D)=0,\qquad
 H^1(K_D)=(S^w)^{n+1}/\operatorname{diag}S^w,
 \qquad H^2(K_D)=H^2(D_n^{w,*}).                         \tag{136.DR8}
\]
Use the coefficient cone
\[
 P_N(K)^i=K^i\oplus K^{i+1},\qquad
            d(x,y)=(dx-Ny,-dy).
\]
Its usual exact sequence, valid without flatness, now becomes
\[
 0\longrightarrow Q_{n,w,N}\longrightarrow H^1(P_N(K_A))
       \longrightarrow H^2(A_n^{w,*})[N]\longrightarrow0.
                                                               \tag{136.DR9}
\]
For a target $C_j$ the rightmost group vanishes, giving the
canonical identification $H^1(P_N(K_{C_j}))=Q_{n,w,N}$.
Each $f_j$ induces a retraction of the inclusion of $Q_{n,w,N}$
in (136.DR9), because it is the identity on the coefficients.
The difference of two retractions is zero on that subgroup,
and so factors uniquely through the rightmost group. This is
the map (136.DR5).

The pair $(b,z)$ is a degree-one cycle of $P_N(K_A)$ and maps
to $[z]$ in (136.DR9). In the full target cone it has primitive
$(-t_j,-c_j)$ by (136.DR3). Its vertex restriction is
$(-t_j|_{V_n},0)$, so the connecting homomorphism of the relative
sequence identifies its class with exactly (136.DR4).
This proves (136.DR6), including its independence of $b$.
\end{proof}

\paragraph{ARCH-136-FIRST-SKELETON-RESIDUE.}
If all the target maps agree in simplicial degree one under the
common first-skeleton identifications supplied by
AQ-136-FIRST-SKELETON, then all the classes (136.DR4) agree.
Consequently every individual sufficiently high-weight degree-two
Tate pair has a simultaneous lift with common vertex coefficients.

\begin{proof}
The quotient (136.DR2) is identified with $(S^w/NS^w)^n$ by
taking differences between vertex $i$ and vertex zero. Its
$i$th coordinate is obtained by restriction to the edge $[0,i]$.
On that edge the data $z,b$ and the target differential algebra
are identical in every target: they lie in the common $F_1$.
Independence of the choices, already proved above, therefore
gives the same residue on that edge in every target. All $n$
coordinates agree. The existence assertion follows from the
criterion just proved.
\end{proof}

\paragraph{The remaining choice is substantive.}
The free degree-two concentration theorem does not assert that
its maps agree in simplex one. Preserving that agreement during
all higher repairs introduces relative primitives over the whole
first skeleton, whose graph cohomology has additional degree-two
terms. Thus the last corollary supplies a useful sufficient
condition for the individual repair, without silently supplying
its preservation throughout the construction. The inherited
AQ-136-HIGHER-COMPLETION example has unequal specified residue
classes and remains a counterexample to arbitrary prescribed-map
completion; neither theorem here contradicts it.
